% ============================================================
%  F. C. Sibbern, Udaf Gabrielis' Breve til og fra Hjemmet.
%    Samlet og udgivet af Frederik Christian Sibbern.
%    Med et Forord af Prof. Harald Høffding. Femte Udgave.
%  Kjøbenhavn: Det Reitzelske Forlag (George C. Grøn);
%    Fr. Bagges Bogtrykkeri, 1893.
%  TRANSCRIPTION (Danish)
%  Source: Google Books scan (Harvard College Library, Widener copy),
%    ~/bibliotek/Sibbern, Frederik/Udaf_Gabrielis_breve_til_og_fra_hjemmet.pdf
%  364 PDF pages. Roman/antiqua throughout (NOT Fraktur).
%  PAGE MAP (see RESUME-NOTES.md for full detail):
%    Two independent roman front-matter sequences:
%      Høffding's Forord (roman): printed = PDF - 10  (p.III = PDF 13).
%      Sibbern's Tilegnelse (roman): printed = PDF - 46 (p.V = PDF 51).
%    Body / the letters (arabic): printed = PDF - 54  (p.1 = PDF 55).
%    Body ends ~p.302 (PDF ~356); tail leaves are library slips.
%  BATCH 1 (scaffold): preamble, title page, and Sibbern's Tilegnelse
%    (dedication to Bishop Mynster, pp. V--VIII) — image-verified.
%    Høffding's Forord (pp. III--XXXVII) and the body of letters
%    (pp. 1--302) are NOT yet transcribed.
% ============================================================
\documentclass[12pt,a4paper]{book}
\usepackage[utf8]{inputenc}
\usepackage[T1]{fontenc}
\usepackage{libertinus}
\usepackage{libertinust1math}
\usepackage{textalpha}
\usepackage[danish]{babel}
\usepackage[protrusion=true,expansion=true]{microtype}
\usepackage{setspace}
\usepackage[margin=1.2in]{geometry}
\usepackage{fancyhdr}
\setlength{\headheight}{28pt}
\addtolength{\topmargin}{-13.5pt}
\usepackage{hyperref}

\hypersetup{
  pdftitle={Udaf Gabrielis' Breve til og fra Hjemmet},
  pdfauthor={F. C. Sibbern},
  hidelinks
}

\pagestyle{fancy}
\fancyhf{}
\fancyhead[LE,RO]{\thepage}
\fancyhead[CE]{\textit{\leftmark}}
\fancyhead[CO]{\textit{\rightmark}}

\setstretch{1.3}

\title{\Huge\textbf{Udaf}\\[6pt]
  \Huge\textbf{Gabrielis' Breve}\\[8pt]
  \large \textbf{til og fra Hjemmet.}\\[14pt]
  \normalsize Samlet og udgivet af\\[4pt]
  \large Frederik Christian Sibbern.\\[10pt]
  \normalsize Med et Forord af Prof.\ Harald Høffding.\\[10pt]
  \normalsize Femte Udgave.}
\author{}
\date{Kjøbenhavn 1893.\\
  Det Reitzelske Forlag (George C.\ Grøn).\\
  Fr.\ Bagges Bogtrykkeri.}

\begin{document}
\maketitle
\thispagestyle{empty}
\cleardoublepage

%% ============================================================
%%  FORORD  (Harald Høffding; printed pp. III--XXXVII; PDF 13--47)
%%  NOT yet transcribed. Roman/antiqua, set largely in italic in the
%%  original. Offset: printed = PDF - 10.
%% ============================================================
\chapter*{Forord.}
\markboth{Forord}{Forord}

\begin{center}\textit{[Forord af Harald Høffding, pp.~III--XXXVII --- ikke transskriberet endnu.]}\end{center}

\cleardoublepage

%% ============================================================
%%  TILEGNELSE  (Sibbern's dedication to Bishop J. P. Mynster;
%%  printed pp. V--VIII; PDF 51--54). Dated 4 Oct. 1850, with a
%%  note to the 4th edition dated July 1870. Image-verified.
%%  Offset for this sequence: printed = PDF - 46.
%% ============================================================
\chapter*{Til Deres Excellence Biskop J.\ P.\ Mynster.}
\markboth{Tilegnelse}{Tilegnelse}

\begin{center}\rule{0.28\linewidth}{0.4pt}\end{center}

„Jeg tænkte paa de Dage fra fordum og de Aar, som ere henrundne“ ---
hvor ofte ere ikke disse Assaphs Ord vendte tilbage i min Erindring,
fra den Tid af, da jeg i Aaret 1812 paa Aarets sidste Dag i
Aftendæmringen hørte en Prædiken over dem i den store Domkirke i
Breslau, hvor jeg dengang opholdt mig, for at være i Nærheden af
Steffens. Men især vendte jeg ofte tilbage til disse Ord fra den Tid
af, da jeg begyndte at erholde og samle de Breve af Gabrielis, som jeg
her meddeler, og tænkte paa, at, naar jeg en Gang kom til at udgive
dem, maatte jeg og vilde jeg lade dem ledsages af en Tilegnelse til
Dem. Det er nu 25 Aar siden, at jeg sendte Dem „Gabrielis' efterladte
Breve“ i Haandskrivten, og fra Dem med stor Glæde modtog den Recension
deraf, som blev trykket foran dem. Da vare Tiderne ganske andre, end
de nu ere, men ogsaa da stredes og kæmpedes der i Gemyterne og af
Gemyterne i Forvirring og Uro; det var i kirkelig Henseende, at der
blev stridt, og jeg udgav netop hine Breve (som da allerede, paa nogle
faa nær, havde ligget hos mig i over 9 Aar), fordi jeg forestillede
mig, at de maaskee kunde virke til, at Tankerne hos os i de offentlige
Omtalers Gebeet, kunde vende sig fra Spørgsmaalet om de rette første
Kilder til Livets helligste Vande, hen til Omsorgen for den rette
Gjenfødelse og Qvægelse ved disse Vande selv. Hine Breve lade os see
Gabrielis i en Krise, og det er kun denne Krise, de føre frem for os,
ladende Enhver i egne Tanker forfølge den til dens Opløsning eller
Klarelse; men Tilværelsens hellige Vande rinde og raade deri. Jeg
havde efter hiin Tid ofte det Haab, engang at skulle kunne knytte
tvende Samlinger af Breve til hiin første, og at deraf den ene skulde
kunne kaldes „Gabrielis' Breve til Hjemmet“, den anden, hans „Breve fra
Hjemmet“. Men kun fra nogle faa Uger, men vistnok for ham indholdsrige
og betydningsfulde Uger, fra en Tid, som falder 10 Aar efter hiin Tid i
hans Liv, har jeg kunnet komme i Besiddelse af Breve fra hans Haand, og
dem meddeler jeg nu.

„Jeg tænkte paa de Dage fra fordum og de Aar, som ere henrundne“: den
hellige Sanger udbryder i disse Ord i sin Trængsels Tid, paa sin
Beængstelses Dag; men idet han mindes sit Folks underfulde Førelser,
vende hans Tanker sig snart til Herrens Gjerninger i fordums Dage,
Herrens Idrætter og Veie, og hans Tilbageblik bliver hans store
Fortrøstning. I Sandhed det forekommer mig, at vi her i Danmark have
megen Aarsag til at skue tilbage med lignende Følelser.

„Jeg tænkte paa de Dage af fordum og paa de Aar af Evighederne“, saa
lyde bogstaveligen Ordene i Grundsproget, og saaledes giver vor
kirkelige Oversættelse os dem. „De Tider af Evighederne“, hvor ganske
anderledes gribe saadanne Ord ind i os ved den Betydning, vi saa
naturligen føle os drevne til at lægge ind i dem, end naar vi blot
holde os til den Betydning, vi maae tillægge dem hos den hellige
Sanger.

De Timer, som leves af Evigheden, leves ogsaa for Evigheden, og leves
tillige for lange, fjerne, jordiske Tider, som skulle komme. Herren har
udkaaret Dem til at være hans Røst til Udbredelse af mange saadanne
Virkninger i vort Fædreneland, og han lader Dem endnu i høi Alder
vedblive i saadan velgjørende Virken. Jeg erfarer det idelig: Deres Ord
og Tale har været til stor Velsignelse, som i min nærmeste Kreds, saa i
mange Kredse, og Gud bevare os Dem endnu længe, og bevare Dem Aandens
Kraft, indtil de Dage komme, da De selv villigen og med Tilfredshed vil
sige: Herre! nu lade Du Din Tjener vare i Fred!

\begin{flushright}
Den 4de Oktober 1850.\\
\textit{Sibbern.}
\end{flushright}

\begin{center}\rule{0.18\linewidth}{0.4pt}\end{center}

Ved nærværende fjerde, paa nogle enkelte Ord nær, aldeles uforandrede
Udgave, bør jeg ikke undlade (hvad der allerede tidligere burde skeet)
at gjøre opmærksom paa, at Gabrielis ved Zodiakallyset forstaaer den
fra Vest til Nord og derfra til Øst fremskridende Lysning, som vi takke
for vore saakaldte lyse Nætter. Astronomerne forstaae ved det nævnte
Udtryk noget ganske Andet.

\begin{flushright}
Juli 1870.\\
\textit{Sibbern.}
\end{flushright}

\cleardoublepage

%% ============================================================
%%  BODY --- The letters (arabic pagination; printed pp. 1--~302;
%%  PDF 55--~356). Offset: printed = PDF - 54.
%%  Epistolary; letters dated. Divided later into parts ("til
%%  Hjemmet" / "fra Hjemmet" per Sibbern's Tilegnelse).
%%  Transcribed & image-verified so far: pp. 1--10 (PDF 55--64).
%% ============================================================
\cleardoublepage
\markboth{Gabrielis' Breve}{Gabrielis' Breve}

% [printed p.1 / PDF 55]
\begin{center}
Fredag den 21de Mai 1824. Aften.
\end{center}
\medskip

Jeg fornyer engang igjen forrige Dages Minder og Følelser. Efter en
lang Fodtour sidder jeg nu her i et eensomt Gjæsteværelse, for at skrive
til Dig. Hvad der Dagen igjennem har bevæget sig i mig, skal i en mild
Efterklang udgydes paa Papiret, og Sjælen skal imidlertid nyde sig selv
i sin stille hyggelige Væren hos sig selv. Nylig raabte Vægteren, at
Klokken var ti, og at Vinden var Sydvest.

Jeg gik igaar Eftermiddag hjemme fra. Da jeg saae den blaae Luft
gjennem begge Lydhullerne i Marslevs Kirketaarn, og altsaa efter
Sigende var netop Halvveien til Odense, stod Solen allerede meget lavt.
Den brød just frem bag en Sky og kastede sin rødgule Glands hen over
Egnen. Som jeg gik lidt videre, straalede nogle Huse af Solens Skin paa
Vinduerne; inden jeg naaede til Odense, straalede mange Lys som
stjernede Lyspuncter hen til mig fra Husene,
% [printed p.2 / PDF 56]
frem under de mange lave Straatage. Ogsaa fra den gamle Marens Vinduer
straalede et Lys. Jeg forestillede mig hende, hvorledes hun nu lige sad
og læste i Müllers Huuspostil for sine Børn, og det af det gamle
forslidte Exemplar uden Begyndelse og uden Ende. Jeg forærede hende
ifjor et heelt og nyt Exemplar, som hun vilde gjemme til sin næstyngste
Datter, sagde hun, naar hun til Mortensdags Tider dette Aar skal ud og
tjene for sit Brød. Imidlertid vilde de læse deri hver Gang, de kom til
en Prædiken, som i hendes gamle Bog ikke var at læse, fordi den enten
var reent borte eller Blade deraf vare afrevne. Men i Aften, saa tænkte
jeg ved mig selv, læser hun i den gamle Bog; thi Prædikenen til femte
Søndag efter Paaske, vidste jeg, var heel og ufordærvet; og hun læser
dog helst de Blade, hun nu i saa mange Aar har elsket. At hun ikke, som
forhen, læser den forestaaende Søndags Prædiken om Løverdagen men om
Fredagen, er efter mit af hende strax fulgte Raad, for at den læste
Prædiken ikke skal næsten gaae lige foran for den, som om Søndagen skal
høres. Det er saa godt at gaae med en Prædiken i Sjælen om Løverdagen,
sagde jeg til hende; det slog hende især. Jeg tænkte paa, hvorledes
ogsaa hun hører med til de herlige Mennesker, med hvis Tilværelse vor
Herre har beriget mit
% [printed p.3 / PDF 57]
Liv. Jeg har dog ei saa faa saadanne lige fra Huusmandshusene lige til
Palaierne.

I Odense besøgde jeg intet Menneske. Jeg lod mig i Værtshuset give mit
sædvanlige Værelse og sad i Eenrum. Og hvormed levede jeg? Ikke med mit
lille nye Testament, ikke med det 13de Bind af Goethe. Begge havde jeg
med; ved begges Ord havde jeg paa Veien frydet mit Sind. De to Vers i
Goethes: „die Geheimnisse“, som begynde med den Linie: „Wenn einen
Menschen die Natur erhoben“, have i mange Aar hørt til mine bedste
Klenodier. Nu fik jeg dem dog paa denne Vei lært ret fuldkomment
udenad, saa at de ei skulle slippe mig mere. Med mit nye Testament gik
det mig atter, som oftere. Er det ikke dog besynderligt? Man læser det
hele Brev til Galaterne igjennem, endda i en opvakt Stund. Saa en anden
Gang slaaer man op paa en Spadseretour og læser kun et Par Vers af det
samme Brev, og disse Par Vers fylde og qvæge og nære mere, end forhen
det hele Brev, hvoraf de dog kun udgjøre en lille Deel. Saaledes mætter
til rette Tid lidt Fisk og lidt Brød mere, end det Mere end Dobbelte
ellers.

Men nu --- da jeg nu sad i mit sædvanlige Gjæsteværelse med dets
Knaphed, --- et Bord, to Stole, et Slags Servante med Vandkande og
% [printed p.4 / PDF 58]
Fad, det hvide Haandklæde paa den kridhvide Væg, Alt tæt op til Sengen,
to Spedelys foran mig, thi med eet kan jeg ikke lade mig nøie, --- nu da
jeg sad ved mit Smørrebrød, halvt med Kommenost, halvt med mager Skinke
paa, en Flaske Øl derhos, nu læste jeg hverken i den ene eller den anden
af mine medbragte Bøger, skjønt Alt endnu dæmrede aftenligt i mig; nei,
jeg tog til Berlingsaviserne, som Pigen havde lagt hen til mig. Der gaae
stundom Uger, ja Maaneder hen, da jeg ingen Aviser læser; saa kan jeg da
engang igjen ret gotte mig ved dem i en ledig Stund; dog, Gottelse er
det just ikke; men jeg strakte mig hen i dem, som man efter en trættende
Tour strækker sig hen paa en Sopha eller en Løibænk; og endnu i Sengen
læste jeg min Rest. Og ogsaa herved følte jeg mig ret vel tilmode.
Aanden kommer ei, naar vi kalde, men naar den kalder os. Dog jo, den
kommer ogsaa, naar vi kalde. Men jeg kaldte ei; den havde rørt sig i mig
og varmet mig og sjunget for mig den hele Dag. Nu slumrede den, jeg kan
sige sødeligen, og min Sjæl legede selv slumrende ved dens Barm med
ingen Ting. Hvorfor skulde den ei det?

Idag gik jeg, uden i Odense at have besøgt noget Menneske, igjennem det
deilige sydlige Fyen. Ikke dog, som forrige Aar mere end
% [printed p.5 / PDF 59]
fjorten Dage tidligere, under Kirkeklokkernes høitidelige Klang, hen
mellem pyntede Bønderkoner og Bønderpiger, vevre og flinke, som deres
Mænd og Ungkarle. Dengang var det Søndag; nu var Løverdagen kjendelig
overalt paa Mark og Sti, i Huse og Gaarde. Jeg gik over Brahetrolleborg.
Hvilken Herlighed! Touren gjorde jeg langsomt og i al Mag; ofte hvilede
jeg i disse foraarsfriske Skove; --- Jordbund og Grøvter frodigt grønne
og Løvet ungt og lysegrønt, temmeligt vidt udsprunget overalt; hvide
Anemoner i den grønne Bund; mange gule Blomster stode som Stjerner.
Sommerfugle tumlede sig, utallige Fuglestemmer qviddrede, og med dem
rislede de smaae ilende Væld, der fremkaldes om Foraaret og om Sommeren
forsvinde. Ingen Hovedlandevei fulgte jeg, og agtede det ei, at min Vei
blev en halv Miil længere. Jeg havde hele Dagen for mig, og lod den
jevnt gaae med. Nu og da var jeg inde i et Bondehuus eller en
Bondegaard, og lod dem fortælle mig, hvilken rar Mand deres Præst var,
og hvordan det havde fornøiet dem, da Kongen havde gjort hans Kone til
Frue og hans Datter til Frøken, skjønt han ikke havde skjøttet om det.

Saa kom jeg da til denne By, efterat Mørket var begyndt at falde paa.
Men her er mødt
% [printed p.6 / PDF 60]
mig noget høist Uhyggeligt og Truende; vidste jeg kun, hvad det er.

Jeg troer, Du veed, hvad det er, der hvert Aar paa denne Tid kalder mig
herhen. Jeg skal med min gamle Toldcassererenke i Hjertelighed og
Veltilfredshed mindes og feire den hende saa kjære, saa elskelige, saa
dyrebare Dødsdag. Det er den Dag, da hendes Søsterdatter for fem Aar
siden døede, næsten i samme Time, som, fire Mile fra hende, hendes
forhenværende Elskede. At miste kjære, elskede Sjæle ved Døden, siger
man med Rette, danner en skjøn, en inderlig Forbindelse imellem os og
Himmelen; og saa besynderlig lever min gamle hjertensgode, brave,
forstandige Enkefrue (thi hendes Mand havde en Rang) med Himmelen,
medens hun sidder eller sysler i sin lille Hjørnestue til Haugen, at det
er, som laae Himlen paa den anden Side af de langt opadstrakte
Haugebede, og som om hun hilsede derover hver Dag; men dog pudsler hun
ganske gjærne, hvor hun er, og lader sine Tanker gaae ved sin Rok. Vi
have Søndag i Morgen; paa Tirsdag er Dødsdagen.

Men hvad monne nu her være skeet i denne By? Jeg gyser ved den Tanke, at
jeg i Morgen skal faae Vished om Det, jeg nu kun aner og frygter: at
Handelskrisernes Uveir med Hagel og Lyn er slaaet ned i denne fredelige
Egn, og
% [printed p.7 / PDF 61]
at Ødelæggelsen vil vise sig paa Alles Aasyn. Om Bankerotter har man jo
hørt meget tale. Sikkert have disse ødelagt Kjøbmændene her i Byen, og
da den hele By næsten er indbyrdes beslægtet og besvogret, ere vel Alle
i Sorg og Qvide, kan jeg troe.

Kun min gode Moses Aaron har sikkert holdt sig, og kan maaskee tage
megen af den Handel op, der gaaer til Grunde for de Andre. Hvis saa er,
saa kunde jeg næsten begynde at hade ham for hans, ja jeg tør sige
skadefroe Aasyn i Aften. Han stod for sin Dør, som han pleier; stundom
sætter jeg mig da hos ham paa Træbænken. I Aften lod jeg det dog være;
thi jeg var baade træt og noget øm i Fødderne, og saa gjør man vel i,
især naar man er paa Brostenene, at fortsætte sin Vei uafbrudt, til man
er ved sit endelige Hvilested. Idetmindste skal man ikke sætte sig;
dvæler man, saa dvæle man staaende. Jeg blev da staaende en lille Stund
hos Manden.

Han er en af de halvgamle Jøder, som jeg holder meget af og ved
Leilighed gjærne giver mig af med. Har det nogensinde truffet sig for
Dig, at Du har hørt en saadan halvgammel Jøde af det ægte Slags læse
Noget af det gamle Testament? Denne hebraiske og orientalske Ild i deres
Udtalelse: hvorledes de sige deres: „Schpricht der Herr“ med en
Eftertrykkelighed,
% [printed p.8 / PDF 62]
--- nei man maa have hørt det selv, for at kunne føle, at sit gamle
Testament skulde man, istedetfor at læse det selv, lade sig læse for af
en saadan Jøde. Min gode Moses Aaron har leilighedsviis oftere læst Et
eller Andet af Propheterne for mig.

Han stod da for sin Dør, og jeg standsede hos ham, og spurgte ham, om
Alt stod sig godt og vel til i Byen, og især, om min Enkefrue var vel.
„Ach leider Gottes, nein! Und der ganze Ort, ja da werden Sie Wunder
Gottes høren, --- Alles ist bankerot. Nein Du mein Himmel --- das sind
Christen, das wollen Christen seyn.“ Han sagde Dette med en besynderlig
triumpherende og besynderlig spottende og ironiserende Tone. Mig stødte
det; dobbelt, fordi Glæde over Fordeel ved Andens Skade dog ellers ikke
ligner min gode Moses Aaron, hvem jeg altid har holdt af, som en inderlig
honnet Natur, og tillige som, jeg vil ikke sige en from eller
gudfrygtig, men dog religiøs Mand, der holder fast ved Abrahams, Isaaks
og Jakobs Gud. „Naa, svarede jeg, saa staaer det vel saare godt til i
Deres Stue, og har De været en holden Mand før, saa bliver De det vel
dobbelt nu.“ --- „Ja, deri har De Ret“, sagde han med et, jeg veed ei,
om jeg skal sige Smiil eller Griin. Vi have det saa med hinanden, at jeg
bestandig
% [printed p.9 / PDF 63]
taler Dansk med ham, og at han ogsaa i Almindelighed taler Dansk igjen,
men derimod Tydsk, naar Ilden kommer op i ham. „Men, føiede han til,
Stuesager er det ikke, men Kammersager, reine Kammersachen, Gott strafe
mich, reine Kammersachen. In der Stube treiben die Spinnen ihr Werk. Nun
Sie werdens erfahren; dann wollen wir mit einander reden.“

Jeg gik og tænkte: hvad kan det være? Du maa, for at forstaae hans Ord,
vide, at han boer alene i en Stueetage med en Stue, to Kamre og et
Kjøkken, hvor han selv gjør Alt, endog at hente Vandet fra Pumpen strax
ved i Gaarden. I Stuen har han sin Handel. I Kamret til Gaden derimod ---
i det til Gaarden sover han med en lille Bjæffer af Hund, et aarvaagent
og meget støiende Dyr, næsten utaaleligt, naar en Fremmed kommer, men
just derfor har han Hunden til sin Sikkerhed, --- i Kamret til Gaden,
siger jeg, har han sin Bibel, en stor Foliant, næsten altid opslaaet, som
Fundament for det Øvrige, og saa sin Mendelsohn, som han er stolt af;
han gjør sig ogsaa til af Spinoza, som han har et Portrait af; men en
Oversættelse af Spinoza, han har, vil ei smage ham. Desuden har han
mangen anden Bog. At tale om Gud og om Verdensforholdene med mig, naar
jeg kommer, holder han meget af, og jeg besøger
% [printed p.10 / PDF 64]
ham gjærne. Saa sidde vi altid i Kamret; og nu er det blevet til et Ord
imellem os, at kalde hans Handelsaffairer Stuesagerne, men vore lærde og
curiøse Forhandlinger kalde vi Kammersagerne. Men jeg frygter dog, at de
Bankerotter, man har hørt tale om i denne Tid, have slaaet ned her og
ødelagt.

Jeg traadte paa Veien ind i en Boutik, for, medens jeg huskede det, at
kjøbe mig en Blyant, som jeg havde glemt at kjøbe mig i Odense, hvor jeg
havde tænkt at gjøre det, og som jeg paa min hele Tour havde følt Savnet
af. Jeg er vel bekjendt i Boutiken som i hele Byen. Overhovedet holder
jeg af, naar jeg skal have Noget kjøbt eller skal have Noget bestilt hos
en Haandværker, at gaae sligt Ærinde selv. Det fornøier mig at komme ind
i Handelsmandens, i Haandværkerens hele indre Tilværelse, at see ham i
hans Dont og dvæle hos ham. Tæt til den Boutik, jeg gik ind i, støder et
Kammer, hvor Manden pleier at staae ved sin Pult. Naar jeg kommer, gaaer
det aldrig af uden en Deel Passiar, først med Boutikdrengen, saa med
Manden selv, som gjærne strax kommer. Saa maae de vise mig, hvad de have
faaet af Nyt, siden jeg var her sidst, og har Drengen en ny Trøie af
anden Farve, end den forrige, saa glemmer jeg ikke at raisonnere
derover, og saadan gaaer
% [Brev fortsætter --- printed p.11 (PDF 65) og frem: IKKE transskriberet endnu.]

\end{document}
